\section{Conclusion}
\label{sec:conclusion}

This report presented \textbf{MedRAGAgents}, a production-grade multi-agent architecture designed for complex biomedical reasoning on the MedQA-USMLE benchmark. By decomposing single monolithic prompts into specialized agent roles—Domain Classification (\texttt{DomainAgent}), Parallel Specialty Analysis (\texttt{AnalysisAgent}), Diagnostic Synthesis (\texttt{SynthesisAgent}), and Consensus Verification (\texttt{VerifierAgent})—the system addresses the cognitive overload and hallucination risks common in direct baseline models.

Integrating biomedical dense vector retrieval via MedCPT over a preprocessed textbook corpus supplies precise factual evidence at the synthesis stage, mitigating knowledge cutoff limitations. The dual-layer memory infrastructure separates transient execution states in an in-RAM scratchpad (\texttt{ShortTermMemory}) from persistent disk records (\texttt{LongTermMemory}) indexed by canonical SHA-256 digests, guaranteeing zero-latency replay on previously solved vignettes while preventing cross-question context contamination.

Empirical evaluation across progressive variants (V0 through V3) on the MedQA-USMLE test set demonstrated measurable accuracy gains over direct chain-of-thought baselines while maintaining an output formatting invalidity rate of $0.0\%$. Non-parametric bootstrap confidence intervals and paired McNemar test analyses confirm the statistical reliability of the multi-agent approach. 

The complete codebase, TikZ architectural specifications, evaluation scripts, and serialized prediction outputs provide a foundation for future research in multi-specialty clinical decision support systems.
